% Captions for harmonic-scattering-loop-coupling panels.
% Include via % Captions for harmonic-scattering-loop-coupling panels.
% Include via % Captions for harmonic-scattering-loop-coupling panels.
% Include via % Captions for harmonic-scattering-loop-coupling panels.
% Include via \input{figures/loop-captions.tex}.

\begin{figure*}[!htbp]
    \centering
    \includegraphics[width=\textwidth]{figures/panel_1_graph.png}
    \caption{\textbf{Molecular harmonic graph and its cycle structure.}
    (\textbf{A})~Stem plot of benzene's five lowest IR-active vibrational modes at $\omega\in\{673,1038,1486,3068,3099\}$~cm$^{-1}$ (NIST CCCBDB). Each mode is a vertex of the harmonic graph.
    (\textbf{B})~Harmonic edges plotted by characteristic frequency $(p\omega_i+q\omega_j)/(p+q)$ versus fractional mistuning $\delta$ (log ordinate), for $(\eta_{\max},\delta_{\mathrm{tol}})=(10,0.05)$. Seven harmonic edges form over five modes, giving cycle rank $C=|E_h|-|V|+K=7-5+1=3$.
    (\textbf{C})~Median cycle rank as a function of molecular size (modes) for random log-uniform frequency draws in $[500,3200]$~cm$^{-1}$, $N\in[4,12]$, 40 realisations per $N$. Shaded band: interquartile range. Cycle rank grows approximately quadratically in $N$: larger polyatomics reach higher $C$ and thus higher multi-source capacity.
    (\textbf{D})~Three-dimensional scatter of benzene's harmonic edges in the $(\omega_i,\omega_j,\delta)$ space, coloured by $\log_{10}\delta$. Edges with the tightest mistuning (dark points) dominate coherent loop traversal; edges with larger $\delta$ contribute fewer loop iterations before dephasing.}
    \label{fig:panel1_graph}
\end{figure*}


\begin{figure*}[!htbp]
    \centering
    \includegraphics[width=\textwidth]{figures/panel_2_coupling.png}
    \caption{\textbf{The loop-coupling map assigning sources to harmonic edges.}
    (\textbf{A})~Frequency-selectivity curves for each of benzene's seven harmonic edges. Each edge is a Lorentzian peaked at its characteristic frequency with width set by its own mistuning $\delta$. Narrow, well-separated peaks imply that distinct source wavelengths couple to distinct edges — the map is near-injective on wavelength.
    (\textbf{B})~Angular-selectivity curves $|\hat n\cdot(\mu_i\times\mu_j)|$ for five edges as a function of source azimuth in the equatorial plane. Different transition dipole orientations produce orthogonal selection patterns on the sphere, providing the second independent channel of the coupling map (direction) alongside wavelength.
    (\textbf{C})~Combined coupling (frequency selectivity times angular selectivity) along a wavelength sweep at fixed azimuth, one curve per edge. Peaks at different wavelengths demonstrate that sources at distinct wavelengths populate disjoint subsets of the edge set.
    (\textbf{D})~Three-dimensional surface of the coupling map restricted to the tightest-mistuning edge: coupling amplitude as a function of source wavelength $\lambda$ and azimuth. The peak is sharp in $\lambda$ and structured in azimuth; the joint parametrisation is the explicit loop-coupling map $\Phi$ from source attributes to harmonic-edge index.}
    \label{fig:panel2_coupling}
\end{figure*}


\begin{figure*}[!htbp]
    \centering
    \includegraphics[width=\textwidth]{figures/panel_3_reconstruction.png}
    \caption{\textbf{Multi-source reconstruction through the benzene-like harmonic resonator ($C=3$, $K=4$ sources, $\kappa(A)=104$).}
    (\textbf{A})~Reconstruction error $\|\Svec-\hat{\Svec}\|_2/\|\Svec\|_2$ versus noise amplitude $\sigma$ on log--log axes (black), compared to the condition-number bound $\kappa\sigma$ (red dashed). At zero noise the error reaches double-precision machine epsilon ($2.8\times 10^{-15}$); at finite noise the error tracks the $\kappa\sigma$ prediction to within a Tikhonov-regularisation bias term.
    (\textbf{B})~Recovered versus true source amplitudes (real part) at $\sigma=10^{-3}$: grouped bars show amplitude of each of the four sources, with recovered values within 29\% of ground truth. At lower noise the agreement is exact.
    (\textbf{C})~Singular values of the transfer matrix $A$ on semilog ordinate. All four singular values are non-zero and span a single factor of $10^{2}$, confirming the full-rank condition of Theorem~\ref{thm:rank} and setting the condition number $\kappa(A)=104$.
    (\textbf{D})~Three-dimensional error surface over $(\log_{10}\sigma,\mathrm{trial\ index})$ showing reproducibility across ten independent noise realisations. The surface's shallow ridge along increasing $\sigma$ and flat dependence on trial index indicates that the reconstruction is stable under noise-sampling variations; the dominant error is noise amplification, not reconstruction instability.}
    \label{fig:panel3_reconstruction}
\end{figure*}


\begin{figure*}[!htbp]
    \centering
    \includegraphics[width=\textwidth]{figures/panel_4_viscosity.png}
    \caption{\textbf{Viscosity--refractive-index relation across twelve liquids at $20^\circ$C.}
    (\textbf{A})~$\mu$ versus $n_r$ (log ordinate) for twelve liquids, coloured by chemical class: nonpolar (blue), polar-aprotic (green), aromatic (purple), hydrogen-bonding (red). A class-dependent $(\mu,n_r)$ trend is visible: within each class the scaling is monotone, but the proportionality constant differs by class.
    (\textbf{B})~Hydrogen-bonding subset (seven liquids: methanol, water, ethanol, ethylene glycol, olive oil, glycerol, castor oil) on log ordinate. Labels adjacent to each point. Spearman rank correlation $\rho_s=+0.964$ with $p=4.5\times 10^{-4}$, Kendall concordance $20/21=95\%$: the within-class monotonicity is clear and statistically significant.
    (\textbf{C})~Per-liquid concordance rate bar chart: for each liquid, fraction of all pairs with every other liquid that respect the $\mu$--$n_r$ ordering. Bars coloured by chemical class. Hydrogen-bonding liquids achieve concordance above 70\%; nonpolar and polar-aprotic liquids sit near 50\% (uncorrelated across classes). The grey dashed line marks the chance level of $0.5$.
    (\textbf{D})~Three-dimensional scatter of $(n_r,\log_{10}\mu,\mathrm{liquid\ index})$ coloured by chemical class. The data cluster by class, with the H-bond class (red) forming the steep monotone branch from water to castor oil and glycerol.}
    \label{fig:panel4_viscosity}
\end{figure*}


\begin{figure*}[!htbp]
    \centering
    \includegraphics[width=\textwidth]{figures/panel_5_conditioning.png}
    \caption{\textbf{Condition number of the transfer matrix as a function of cycle rank $C$ (60 random harmonic graphs per $C$, $K=C+1$ sources).}
    (\textbf{A})~Median $\kappa(A)$ on log--log axes (blue circles) with interquartile band and power-law fit $\kappa=\beta C^\alpha$ (red dashed). Fitted exponent $\alpha=1.25$ — well below the quadratic threshold of $\alpha=2$ — confirming that noise amplification grows slower than the capacity $K=C+1$.
    (\textbf{B})~Boxplot of $\kappa(A)$ across 60 random realisations per $C$, on linear ordinate. Medians: 2.8 ($C=1$), 6.1 ($C=2$), 8.6 ($C=3$), 18.6 ($C=5$), 29.8 ($C=7$), 49.1 ($C=10$). Whiskers and outliers demonstrate that the spread is modest at every $C$.
    (\textbf{C})~Tolerable noise amplitude $\sigma_{\mathrm{tol}}=0.1/\kappa(A)$ at which reconstruction error remains below 10\%, plotted versus source count $K=C+1$. Green band: IQR. At $K=11$ (cycle rank ten), the median system still tolerates $\sigma\sim 2\times 10^{-3}$, adequate for many practical signal-to-noise regimes.
    (\textbf{D})~Three-dimensional surface of $\log_{10}\kappa$ over the $(C,\mathrm{percentile})$ plane, percentiles sampled at $10$\%--$90$\%. The surface rises monotonically with $C$ and with percentile; its shallow curvature in $C$ visualises the sub-quadratic growth, while its shallow curvature across percentiles indicates that the worst-case is not catastrophically worse than the median.}
    \label{fig:panel5_conditioning}
\end{figure*}
.

\begin{figure*}[!htbp]
    \centering
    \includegraphics[width=\textwidth]{figures/panel_1_graph.png}
    \caption{\textbf{Molecular harmonic graph and its cycle structure.}
    (\textbf{A})~Stem plot of benzene's five lowest IR-active vibrational modes at $\omega\in\{673,1038,1486,3068,3099\}$~cm$^{-1}$ (NIST CCCBDB). Each mode is a vertex of the harmonic graph.
    (\textbf{B})~Harmonic edges plotted by characteristic frequency $(p\omega_i+q\omega_j)/(p+q)$ versus fractional mistuning $\delta$ (log ordinate), for $(\eta_{\max},\delta_{\mathrm{tol}})=(10,0.05)$. Seven harmonic edges form over five modes, giving cycle rank $C=|E_h|-|V|+K=7-5+1=3$.
    (\textbf{C})~Median cycle rank as a function of molecular size (modes) for random log-uniform frequency draws in $[500,3200]$~cm$^{-1}$, $N\in[4,12]$, 40 realisations per $N$. Shaded band: interquartile range. Cycle rank grows approximately quadratically in $N$: larger polyatomics reach higher $C$ and thus higher multi-source capacity.
    (\textbf{D})~Three-dimensional scatter of benzene's harmonic edges in the $(\omega_i,\omega_j,\delta)$ space, coloured by $\log_{10}\delta$. Edges with the tightest mistuning (dark points) dominate coherent loop traversal; edges with larger $\delta$ contribute fewer loop iterations before dephasing.}
    \label{fig:panel1_graph}
\end{figure*}


\begin{figure*}[!htbp]
    \centering
    \includegraphics[width=\textwidth]{figures/panel_2_coupling.png}
    \caption{\textbf{The loop-coupling map assigning sources to harmonic edges.}
    (\textbf{A})~Frequency-selectivity curves for each of benzene's seven harmonic edges. Each edge is a Lorentzian peaked at its characteristic frequency with width set by its own mistuning $\delta$. Narrow, well-separated peaks imply that distinct source wavelengths couple to distinct edges — the map is near-injective on wavelength.
    (\textbf{B})~Angular-selectivity curves $|\hat n\cdot(\mu_i\times\mu_j)|$ for five edges as a function of source azimuth in the equatorial plane. Different transition dipole orientations produce orthogonal selection patterns on the sphere, providing the second independent channel of the coupling map (direction) alongside wavelength.
    (\textbf{C})~Combined coupling (frequency selectivity times angular selectivity) along a wavelength sweep at fixed azimuth, one curve per edge. Peaks at different wavelengths demonstrate that sources at distinct wavelengths populate disjoint subsets of the edge set.
    (\textbf{D})~Three-dimensional surface of the coupling map restricted to the tightest-mistuning edge: coupling amplitude as a function of source wavelength $\lambda$ and azimuth. The peak is sharp in $\lambda$ and structured in azimuth; the joint parametrisation is the explicit loop-coupling map $\Phi$ from source attributes to harmonic-edge index.}
    \label{fig:panel2_coupling}
\end{figure*}


\begin{figure*}[!htbp]
    \centering
    \includegraphics[width=\textwidth]{figures/panel_3_reconstruction.png}
    \caption{\textbf{Multi-source reconstruction through the benzene-like harmonic resonator ($C=3$, $K=4$ sources, $\kappa(A)=104$).}
    (\textbf{A})~Reconstruction error $\|\Svec-\hat{\Svec}\|_2/\|\Svec\|_2$ versus noise amplitude $\sigma$ on log--log axes (black), compared to the condition-number bound $\kappa\sigma$ (red dashed). At zero noise the error reaches double-precision machine epsilon ($2.8\times 10^{-15}$); at finite noise the error tracks the $\kappa\sigma$ prediction to within a Tikhonov-regularisation bias term.
    (\textbf{B})~Recovered versus true source amplitudes (real part) at $\sigma=10^{-3}$: grouped bars show amplitude of each of the four sources, with recovered values within 29\% of ground truth. At lower noise the agreement is exact.
    (\textbf{C})~Singular values of the transfer matrix $A$ on semilog ordinate. All four singular values are non-zero and span a single factor of $10^{2}$, confirming the full-rank condition of Theorem~\ref{thm:rank} and setting the condition number $\kappa(A)=104$.
    (\textbf{D})~Three-dimensional error surface over $(\log_{10}\sigma,\mathrm{trial\ index})$ showing reproducibility across ten independent noise realisations. The surface's shallow ridge along increasing $\sigma$ and flat dependence on trial index indicates that the reconstruction is stable under noise-sampling variations; the dominant error is noise amplification, not reconstruction instability.}
    \label{fig:panel3_reconstruction}
\end{figure*}


\begin{figure*}[!htbp]
    \centering
    \includegraphics[width=\textwidth]{figures/panel_4_viscosity.png}
    \caption{\textbf{Viscosity--refractive-index relation across twelve liquids at $20^\circ$C.}
    (\textbf{A})~$\mu$ versus $n_r$ (log ordinate) for twelve liquids, coloured by chemical class: nonpolar (blue), polar-aprotic (green), aromatic (purple), hydrogen-bonding (red). A class-dependent $(\mu,n_r)$ trend is visible: within each class the scaling is monotone, but the proportionality constant differs by class.
    (\textbf{B})~Hydrogen-bonding subset (seven liquids: methanol, water, ethanol, ethylene glycol, olive oil, glycerol, castor oil) on log ordinate. Labels adjacent to each point. Spearman rank correlation $\rho_s=+0.964$ with $p=4.5\times 10^{-4}$, Kendall concordance $20/21=95\%$: the within-class monotonicity is clear and statistically significant.
    (\textbf{C})~Per-liquid concordance rate bar chart: for each liquid, fraction of all pairs with every other liquid that respect the $\mu$--$n_r$ ordering. Bars coloured by chemical class. Hydrogen-bonding liquids achieve concordance above 70\%; nonpolar and polar-aprotic liquids sit near 50\% (uncorrelated across classes). The grey dashed line marks the chance level of $0.5$.
    (\textbf{D})~Three-dimensional scatter of $(n_r,\log_{10}\mu,\mathrm{liquid\ index})$ coloured by chemical class. The data cluster by class, with the H-bond class (red) forming the steep monotone branch from water to castor oil and glycerol.}
    \label{fig:panel4_viscosity}
\end{figure*}


\begin{figure*}[!htbp]
    \centering
    \includegraphics[width=\textwidth]{figures/panel_5_conditioning.png}
    \caption{\textbf{Condition number of the transfer matrix as a function of cycle rank $C$ (60 random harmonic graphs per $C$, $K=C+1$ sources).}
    (\textbf{A})~Median $\kappa(A)$ on log--log axes (blue circles) with interquartile band and power-law fit $\kappa=\beta C^\alpha$ (red dashed). Fitted exponent $\alpha=1.25$ — well below the quadratic threshold of $\alpha=2$ — confirming that noise amplification grows slower than the capacity $K=C+1$.
    (\textbf{B})~Boxplot of $\kappa(A)$ across 60 random realisations per $C$, on linear ordinate. Medians: 2.8 ($C=1$), 6.1 ($C=2$), 8.6 ($C=3$), 18.6 ($C=5$), 29.8 ($C=7$), 49.1 ($C=10$). Whiskers and outliers demonstrate that the spread is modest at every $C$.
    (\textbf{C})~Tolerable noise amplitude $\sigma_{\mathrm{tol}}=0.1/\kappa(A)$ at which reconstruction error remains below 10\%, plotted versus source count $K=C+1$. Green band: IQR. At $K=11$ (cycle rank ten), the median system still tolerates $\sigma\sim 2\times 10^{-3}$, adequate for many practical signal-to-noise regimes.
    (\textbf{D})~Three-dimensional surface of $\log_{10}\kappa$ over the $(C,\mathrm{percentile})$ plane, percentiles sampled at $10$\%--$90$\%. The surface rises monotonically with $C$ and with percentile; its shallow curvature in $C$ visualises the sub-quadratic growth, while its shallow curvature across percentiles indicates that the worst-case is not catastrophically worse than the median.}
    \label{fig:panel5_conditioning}
\end{figure*}
.

\begin{figure*}[!htbp]
    \centering
    \includegraphics[width=\textwidth]{figures/panel_1_graph.png}
    \caption{\textbf{Molecular harmonic graph and its cycle structure.}
    (\textbf{A})~Stem plot of benzene's five lowest IR-active vibrational modes at $\omega\in\{673,1038,1486,3068,3099\}$~cm$^{-1}$ (NIST CCCBDB). Each mode is a vertex of the harmonic graph.
    (\textbf{B})~Harmonic edges plotted by characteristic frequency $(p\omega_i+q\omega_j)/(p+q)$ versus fractional mistuning $\delta$ (log ordinate), for $(\eta_{\max},\delta_{\mathrm{tol}})=(10,0.05)$. Seven harmonic edges form over five modes, giving cycle rank $C=|E_h|-|V|+K=7-5+1=3$.
    (\textbf{C})~Median cycle rank as a function of molecular size (modes) for random log-uniform frequency draws in $[500,3200]$~cm$^{-1}$, $N\in[4,12]$, 40 realisations per $N$. Shaded band: interquartile range. Cycle rank grows approximately quadratically in $N$: larger polyatomics reach higher $C$ and thus higher multi-source capacity.
    (\textbf{D})~Three-dimensional scatter of benzene's harmonic edges in the $(\omega_i,\omega_j,\delta)$ space, coloured by $\log_{10}\delta$. Edges with the tightest mistuning (dark points) dominate coherent loop traversal; edges with larger $\delta$ contribute fewer loop iterations before dephasing.}
    \label{fig:panel1_graph}
\end{figure*}


\begin{figure*}[!htbp]
    \centering
    \includegraphics[width=\textwidth]{figures/panel_2_coupling.png}
    \caption{\textbf{The loop-coupling map assigning sources to harmonic edges.}
    (\textbf{A})~Frequency-selectivity curves for each of benzene's seven harmonic edges. Each edge is a Lorentzian peaked at its characteristic frequency with width set by its own mistuning $\delta$. Narrow, well-separated peaks imply that distinct source wavelengths couple to distinct edges — the map is near-injective on wavelength.
    (\textbf{B})~Angular-selectivity curves $|\hat n\cdot(\mu_i\times\mu_j)|$ for five edges as a function of source azimuth in the equatorial plane. Different transition dipole orientations produce orthogonal selection patterns on the sphere, providing the second independent channel of the coupling map (direction) alongside wavelength.
    (\textbf{C})~Combined coupling (frequency selectivity times angular selectivity) along a wavelength sweep at fixed azimuth, one curve per edge. Peaks at different wavelengths demonstrate that sources at distinct wavelengths populate disjoint subsets of the edge set.
    (\textbf{D})~Three-dimensional surface of the coupling map restricted to the tightest-mistuning edge: coupling amplitude as a function of source wavelength $\lambda$ and azimuth. The peak is sharp in $\lambda$ and structured in azimuth; the joint parametrisation is the explicit loop-coupling map $\Phi$ from source attributes to harmonic-edge index.}
    \label{fig:panel2_coupling}
\end{figure*}


\begin{figure*}[!htbp]
    \centering
    \includegraphics[width=\textwidth]{figures/panel_3_reconstruction.png}
    \caption{\textbf{Multi-source reconstruction through the benzene-like harmonic resonator ($C=3$, $K=4$ sources, $\kappa(A)=104$).}
    (\textbf{A})~Reconstruction error $\|\Svec-\hat{\Svec}\|_2/\|\Svec\|_2$ versus noise amplitude $\sigma$ on log--log axes (black), compared to the condition-number bound $\kappa\sigma$ (red dashed). At zero noise the error reaches double-precision machine epsilon ($2.8\times 10^{-15}$); at finite noise the error tracks the $\kappa\sigma$ prediction to within a Tikhonov-regularisation bias term.
    (\textbf{B})~Recovered versus true source amplitudes (real part) at $\sigma=10^{-3}$: grouped bars show amplitude of each of the four sources, with recovered values within 29\% of ground truth. At lower noise the agreement is exact.
    (\textbf{C})~Singular values of the transfer matrix $A$ on semilog ordinate. All four singular values are non-zero and span a single factor of $10^{2}$, confirming the full-rank condition of Theorem~\ref{thm:rank} and setting the condition number $\kappa(A)=104$.
    (\textbf{D})~Three-dimensional error surface over $(\log_{10}\sigma,\mathrm{trial\ index})$ showing reproducibility across ten independent noise realisations. The surface's shallow ridge along increasing $\sigma$ and flat dependence on trial index indicates that the reconstruction is stable under noise-sampling variations; the dominant error is noise amplification, not reconstruction instability.}
    \label{fig:panel3_reconstruction}
\end{figure*}


\begin{figure*}[!htbp]
    \centering
    \includegraphics[width=\textwidth]{figures/panel_4_viscosity.png}
    \caption{\textbf{Viscosity--refractive-index relation across twelve liquids at $20^\circ$C.}
    (\textbf{A})~$\mu$ versus $n_r$ (log ordinate) for twelve liquids, coloured by chemical class: nonpolar (blue), polar-aprotic (green), aromatic (purple), hydrogen-bonding (red). A class-dependent $(\mu,n_r)$ trend is visible: within each class the scaling is monotone, but the proportionality constant differs by class.
    (\textbf{B})~Hydrogen-bonding subset (seven liquids: methanol, water, ethanol, ethylene glycol, olive oil, glycerol, castor oil) on log ordinate. Labels adjacent to each point. Spearman rank correlation $\rho_s=+0.964$ with $p=4.5\times 10^{-4}$, Kendall concordance $20/21=95\%$: the within-class monotonicity is clear and statistically significant.
    (\textbf{C})~Per-liquid concordance rate bar chart: for each liquid, fraction of all pairs with every other liquid that respect the $\mu$--$n_r$ ordering. Bars coloured by chemical class. Hydrogen-bonding liquids achieve concordance above 70\%; nonpolar and polar-aprotic liquids sit near 50\% (uncorrelated across classes). The grey dashed line marks the chance level of $0.5$.
    (\textbf{D})~Three-dimensional scatter of $(n_r,\log_{10}\mu,\mathrm{liquid\ index})$ coloured by chemical class. The data cluster by class, with the H-bond class (red) forming the steep monotone branch from water to castor oil and glycerol.}
    \label{fig:panel4_viscosity}
\end{figure*}


\begin{figure*}[!htbp]
    \centering
    \includegraphics[width=\textwidth]{figures/panel_5_conditioning.png}
    \caption{\textbf{Condition number of the transfer matrix as a function of cycle rank $C$ (60 random harmonic graphs per $C$, $K=C+1$ sources).}
    (\textbf{A})~Median $\kappa(A)$ on log--log axes (blue circles) with interquartile band and power-law fit $\kappa=\beta C^\alpha$ (red dashed). Fitted exponent $\alpha=1.25$ — well below the quadratic threshold of $\alpha=2$ — confirming that noise amplification grows slower than the capacity $K=C+1$.
    (\textbf{B})~Boxplot of $\kappa(A)$ across 60 random realisations per $C$, on linear ordinate. Medians: 2.8 ($C=1$), 6.1 ($C=2$), 8.6 ($C=3$), 18.6 ($C=5$), 29.8 ($C=7$), 49.1 ($C=10$). Whiskers and outliers demonstrate that the spread is modest at every $C$.
    (\textbf{C})~Tolerable noise amplitude $\sigma_{\mathrm{tol}}=0.1/\kappa(A)$ at which reconstruction error remains below 10\%, plotted versus source count $K=C+1$. Green band: IQR. At $K=11$ (cycle rank ten), the median system still tolerates $\sigma\sim 2\times 10^{-3}$, adequate for many practical signal-to-noise regimes.
    (\textbf{D})~Three-dimensional surface of $\log_{10}\kappa$ over the $(C,\mathrm{percentile})$ plane, percentiles sampled at $10$\%--$90$\%. The surface rises monotonically with $C$ and with percentile; its shallow curvature in $C$ visualises the sub-quadratic growth, while its shallow curvature across percentiles indicates that the worst-case is not catastrophically worse than the median.}
    \label{fig:panel5_conditioning}
\end{figure*}
.

\begin{figure*}[!htbp]
    \centering
    \includegraphics[width=\textwidth]{figures/panel_1_graph.png}
    \caption{\textbf{Molecular harmonic graph and its cycle structure.}
    (\textbf{A})~Stem plot of benzene's five lowest IR-active vibrational modes at $\omega\in\{673,1038,1486,3068,3099\}$~cm$^{-1}$ (NIST CCCBDB). Each mode is a vertex of the harmonic graph.
    (\textbf{B})~Harmonic edges plotted by characteristic frequency $(p\omega_i+q\omega_j)/(p+q)$ versus fractional mistuning $\delta$ (log ordinate), for $(\eta_{\max},\delta_{\mathrm{tol}})=(10,0.05)$. Seven harmonic edges form over five modes, giving cycle rank $C=|E_h|-|V|+K=7-5+1=3$.
    (\textbf{C})~Median cycle rank as a function of molecular size (modes) for random log-uniform frequency draws in $[500,3200]$~cm$^{-1}$, $N\in[4,12]$, 40 realisations per $N$. Shaded band: interquartile range. Cycle rank grows approximately quadratically in $N$: larger polyatomics reach higher $C$ and thus higher multi-source capacity.
    (\textbf{D})~Three-dimensional scatter of benzene's harmonic edges in the $(\omega_i,\omega_j,\delta)$ space, coloured by $\log_{10}\delta$. Edges with the tightest mistuning (dark points) dominate coherent loop traversal; edges with larger $\delta$ contribute fewer loop iterations before dephasing.}
    \label{fig:panel1_graph}
\end{figure*}


\begin{figure*}[!htbp]
    \centering
    \includegraphics[width=\textwidth]{figures/panel_2_coupling.png}
    \caption{\textbf{The loop-coupling map assigning sources to harmonic edges.}
    (\textbf{A})~Frequency-selectivity curves for each of benzene's seven harmonic edges. Each edge is a Lorentzian peaked at its characteristic frequency with width set by its own mistuning $\delta$. Narrow, well-separated peaks imply that distinct source wavelengths couple to distinct edges — the map is near-injective on wavelength.
    (\textbf{B})~Angular-selectivity curves $|\hat n\cdot(\mu_i\times\mu_j)|$ for five edges as a function of source azimuth in the equatorial plane. Different transition dipole orientations produce orthogonal selection patterns on the sphere, providing the second independent channel of the coupling map (direction) alongside wavelength.
    (\textbf{C})~Combined coupling (frequency selectivity times angular selectivity) along a wavelength sweep at fixed azimuth, one curve per edge. Peaks at different wavelengths demonstrate that sources at distinct wavelengths populate disjoint subsets of the edge set.
    (\textbf{D})~Three-dimensional surface of the coupling map restricted to the tightest-mistuning edge: coupling amplitude as a function of source wavelength $\lambda$ and azimuth. The peak is sharp in $\lambda$ and structured in azimuth; the joint parametrisation is the explicit loop-coupling map $\Phi$ from source attributes to harmonic-edge index.}
    \label{fig:panel2_coupling}
\end{figure*}


\begin{figure*}[!htbp]
    \centering
    \includegraphics[width=\textwidth]{figures/panel_3_reconstruction.png}
    \caption{\textbf{Multi-source reconstruction through the benzene-like harmonic resonator ($C=3$, $K=4$ sources, $\kappa(A)=104$).}
    (\textbf{A})~Reconstruction error $\|\Svec-\hat{\Svec}\|_2/\|\Svec\|_2$ versus noise amplitude $\sigma$ on log--log axes (black), compared to the condition-number bound $\kappa\sigma$ (red dashed). At zero noise the error reaches double-precision machine epsilon ($2.8\times 10^{-15}$); at finite noise the error tracks the $\kappa\sigma$ prediction to within a Tikhonov-regularisation bias term.
    (\textbf{B})~Recovered versus true source amplitudes (real part) at $\sigma=10^{-3}$: grouped bars show amplitude of each of the four sources, with recovered values within 29\% of ground truth. At lower noise the agreement is exact.
    (\textbf{C})~Singular values of the transfer matrix $A$ on semilog ordinate. All four singular values are non-zero and span a single factor of $10^{2}$, confirming the full-rank condition of Theorem~\ref{thm:rank} and setting the condition number $\kappa(A)=104$.
    (\textbf{D})~Three-dimensional error surface over $(\log_{10}\sigma,\mathrm{trial\ index})$ showing reproducibility across ten independent noise realisations. The surface's shallow ridge along increasing $\sigma$ and flat dependence on trial index indicates that the reconstruction is stable under noise-sampling variations; the dominant error is noise amplification, not reconstruction instability.}
    \label{fig:panel3_reconstruction}
\end{figure*}


\begin{figure*}[!htbp]
    \centering
    \includegraphics[width=\textwidth]{figures/panel_4_viscosity.png}
    \caption{\textbf{Viscosity--refractive-index relation across twelve liquids at $20^\circ$C.}
    (\textbf{A})~$\mu$ versus $n_r$ (log ordinate) for twelve liquids, coloured by chemical class: nonpolar (blue), polar-aprotic (green), aromatic (purple), hydrogen-bonding (red). A class-dependent $(\mu,n_r)$ trend is visible: within each class the scaling is monotone, but the proportionality constant differs by class.
    (\textbf{B})~Hydrogen-bonding subset (seven liquids: methanol, water, ethanol, ethylene glycol, olive oil, glycerol, castor oil) on log ordinate. Labels adjacent to each point. Spearman rank correlation $\rho_s=+0.964$ with $p=4.5\times 10^{-4}$, Kendall concordance $20/21=95\%$: the within-class monotonicity is clear and statistically significant.
    (\textbf{C})~Per-liquid concordance rate bar chart: for each liquid, fraction of all pairs with every other liquid that respect the $\mu$--$n_r$ ordering. Bars coloured by chemical class. Hydrogen-bonding liquids achieve concordance above 70\%; nonpolar and polar-aprotic liquids sit near 50\% (uncorrelated across classes). The grey dashed line marks the chance level of $0.5$.
    (\textbf{D})~Three-dimensional scatter of $(n_r,\log_{10}\mu,\mathrm{liquid\ index})$ coloured by chemical class. The data cluster by class, with the H-bond class (red) forming the steep monotone branch from water to castor oil and glycerol.}
    \label{fig:panel4_viscosity}
\end{figure*}


\begin{figure*}[!htbp]
    \centering
    \includegraphics[width=\textwidth]{figures/panel_5_conditioning.png}
    \caption{\textbf{Condition number of the transfer matrix as a function of cycle rank $C$ (60 random harmonic graphs per $C$, $K=C+1$ sources).}
    (\textbf{A})~Median $\kappa(A)$ on log--log axes (blue circles) with interquartile band and power-law fit $\kappa=\beta C^\alpha$ (red dashed). Fitted exponent $\alpha=1.25$ — well below the quadratic threshold of $\alpha=2$ — confirming that noise amplification grows slower than the capacity $K=C+1$.
    (\textbf{B})~Boxplot of $\kappa(A)$ across 60 random realisations per $C$, on linear ordinate. Medians: 2.8 ($C=1$), 6.1 ($C=2$), 8.6 ($C=3$), 18.6 ($C=5$), 29.8 ($C=7$), 49.1 ($C=10$). Whiskers and outliers demonstrate that the spread is modest at every $C$.
    (\textbf{C})~Tolerable noise amplitude $\sigma_{\mathrm{tol}}=0.1/\kappa(A)$ at which reconstruction error remains below 10\%, plotted versus source count $K=C+1$. Green band: IQR. At $K=11$ (cycle rank ten), the median system still tolerates $\sigma\sim 2\times 10^{-3}$, adequate for many practical signal-to-noise regimes.
    (\textbf{D})~Three-dimensional surface of $\log_{10}\kappa$ over the $(C,\mathrm{percentile})$ plane, percentiles sampled at $10$\%--$90$\%. The surface rises monotonically with $C$ and with percentile; its shallow curvature in $C$ visualises the sub-quadratic growth, while its shallow curvature across percentiles indicates that the worst-case is not catastrophically worse than the median.}
    \label{fig:panel5_conditioning}
\end{figure*}
